% BAB 5 - KESIMPULAN DAN SARAN

\vspace{1cm}
\subsection{Kesimpulan}
\indent
Setelah melalui tahapan perancangan, implementasi, hingga pengujian terhadap Sistem Point of Sale (POS) WellClean Laundry, maka dapat ditarik beberapa kesimpulan utama sebagai jawaban atas pertanyaan penelitian:

\begin{enumerate}
    \item Berdasarkan RQ 1, aplikasi telah berhasil diimplementasikan dengan menggunakan bahasa pemrograman Go dan basis data Supabase. Hal ini memungkinkan transformasi pencatatan konvensional menjadi pengelolaan data operasional laundry yang terpusat, aman, dan dapat diakses secara \textit{real-time}.
    \item Berdasarkan RQ 2, sistem telah berhasil mengintegrasikan fitur automasi verifikasi pembayaran QRIS melalui mekanisme integrasi endpoint notifikasi. Fitur ini memberikan pembaruan signifikan pada proses bisnis dengan memungkinkan validasi transaksi secara otomatis tanpa perlu pengecekan mutasi secara konvensional oleh kasir.
    \item Berdasarkan RQ 3, hasil pengujian unit menggunakan \textit{go tool cover} menunjukkan bahwa fungsi-fungsi krusial seperti keamanan (\textit{hashing password}) dan logika pembayaran telah mencapai tingkat \textit{statement coverage} sebesar 100\%. Hal ini menjamin keandalan sistem dalam menangani proses-proses kritis.
    \item {Berdasarkan RQ 4}, penerapan metode \textit{Scrum} telah membantu proses pengembangan sistem menjadi lebih terstruktur dan adaptif, sehingga seluruh fitur utama mulai dari manajemen pelanggan hingga laporan keuangan dapat diselesaikan sesuai dengan kebutuhan operasional pengguna.
\end{enumerate} 

\vspace{0.5cm}
\subsection{Saran}
\indent
Demi pengembangan sistem yang lebih optimal di masa mendatang, penulis menyarankan beberapa poin peningkatan sebagai berikut:


\begin{enumerate}
    \item Deployment Sistem ke Server Publik: Melakukan publikasi aplikasi ke layanan \textit{cloud hosting} agar sistem dapat diakses secara \textit{online} oleh pemilik dan pelanggan dari mana saja, tidak terbatas pada lingkungan \textit{localhost}.
    \item Pengembangan Aplikasi Mobile Native: Mengembangkan versi aplikasi berbasis Android atau iOS untuk memberikan pengalaman pengguna yang lebih baik dan fitur yang lebih terintegrasi dengan perangkat \textit{smartphone}.
    \item Integrasi Langsung Payment Gateway: Mengembangkan sistem agar dapat terintegrasi langsung dengan API \textit{payment gateway} resmi untuk mengurangi ketergantungan pada aplikasi pihak ketiga dalam menangkap notifikasi pembayaran.
    \item Akselerasi Cakupan Pengujian Unit: Melakukan optimasi pada modul-modul yang saat ini memiliki persentase coverage di bawah 50\% untuk memastikan ketahanan alur logika sistem.
    \item Ekspansi Fitur Notifikasi: Mengintegrasikan layanan WhatsApp Gateway untuk memberikan pembaruan status cucian secara \textit{real-time} kepada pelanggan setelah sistem berhasil di-\textit{deploy}.
\end{enumerate}